\section{从训练最优到推理最优}

\subsection{Chinchilla 比例的局限}

Chinchilla 比例回答的是：\textbf{给定固定训练 FLOP 预算，如何获得最低的训练损失？}但这并非企业真正关心的问题。

\begin{figure}[H]
\centering
\includegraphics[width=0.95\textwidth]{slides-images/slide-046.jpg}
\caption{从训练最优到推理最优：tokens/parameter 比例的演进\protect\footnotemark}
\end{figure}
\footnotetext{Slides 第47页。}

当 Chinchilla 和 Kaplan 的论文发表时（2020--2022），LLM 还不是产品，不需要考虑推理成本。但如今 LLM 已经是创收产品，推理成本是核心关切：

\begin{importantbox}{训练成本是一次性的，推理成本是持续的}
训练一个模型的计算成本只需付一次，但部署后的推理成本随用户使用量\textbf{持续累积}。因此，企业宁愿在训练时多花计算（``过度训练''小模型），也不愿部署一个推理昂贵的大模型。
\end{importantbox}

tokens/parameter 比例的历史演进：
\begin{itemize}
    \item \textbf{GPT-3 (2020)}：约 2 tokens/parameter（严重``训练不足''）
    \item \textbf{Chinchilla (2022)}：约 20 tokens/parameter（训练最优）
    \item \textbf{LLaMA 系列 (2023--)}：100--200 tokens/parameter
    \item \textbf{最新模型 (2024--2025)}：可能超过 1000 tokens/parameter（如 Qwen 用 30T tokens 训练 7B 模型）
\end{itemize}

趋势非常明确：业界正在大幅增加 tokens/parameter 比例，\textbf{用更多的训练计算换取更小、更便宜的推理模型}。

\subsection{Scaling Laws 作为通用工具}

\begin{figure}[H]
\centering
\includegraphics[width=0.95\textwidth]{slides-images/slide-048.jpg}
\caption{Scaling Laws 在文本扩散模型上的应用：同样适用的 IsoFLOP 分析\protect\footnotemark}
\end{figure}
\footnotetext{Slides 第49页。}

Tatsu 以自己学生 Isan 的文本扩散模型研究为例，说明 Scaling Laws 的方法论\textbf{并非局限于自回归 Transformer}。即使在一个全新的生成模型家族上：
\begin{itemize}
    \item IsoFLOP 分析依然产生清晰的 U 型曲线
    \item 最优点依然可靠地遵循幂律 Scaling
    \item 不同模型类型之间仅有常数偏移
\end{itemize}

\begin{figure}[H]
\centering
\includegraphics[width=0.95\textwidth]{slides-images/slide-049.jpg}
\caption{自回归 vs. 扩散模型：两种完全不同的生成模型均展现清晰的 Scaling 行为\protect\footnotemark}
\end{figure}
\footnotetext{Slides 第50页。}

\begin{knowledgebox}{Scaling Laws 的普适性}
Scaling Laws 不是针对特定架构或任务的特殊现象，而是深度学习中的一种\textbf{普遍规律}。只要你的模型足够灵活、数据足够多、训练足够长，就很可能观察到幂律缩放行为。这使得 Scaling Law 的方法论成为\textbf{通用的工程工具}。
\end{knowledgebox}

\subsection{本章小结}

Chinchilla 比例（20 tokens/parameter）只是训练最优的起点。在产品化时代，推理成本的重要性使得 tokens/parameter 比例不断增大。同时，Scaling Laws 的方法论作为通用工具，已经被成功应用到扩散模型等全新领域，展现了其普适性。

%% ========== Part 9: Scaling Law 的实用流程 ==========

